\documentclass[10pt]{extarticle}
\usepackage[a4paper,margin=1.5cm]{geometry}
\usepackage[T1]{fontenc}
\usepackage{amsmath,amssymb}
\usepackage{graphicx}
\usepackage[export]{adjustbox}
\usepackage{multicol}
\usepackage{enumitem}
\usepackage{booktabs,array,tabularx}
\setlength{\columnsep}{0.7cm}
\setlength{\parindent}{0pt}
\setlength{\parskip}{3pt}
\pagestyle{empty}
\makeatletter
% figure registry (figdefs.tex is generated by the extract step)
\newcommand{\qfigdef}[4]{\@namedef{qf@file@#1}{#2}\@namedef{qf@w@#1}{#3}\@namedef{qf@wide@#1}{#4}}
\newcommand{\qf@check}[1]{\@ifundefined{qf@file@#1}{\PackageError{score2}{Unknown figure id #1}{}}{}}
\newcommand{\qf@img}[4][]{\includegraphics[width=#2\dimexpr\@nameuse{qf@w@#4}pt\relax,max width=#3,max height=0.45\textheight,#1]{\@nameuse{qf@file@#4}}}
% Full-width questions: finalize wraps every question that contains a wide figure in \qfullbegin...\qfullend,
% so the WHOLE question (stem, figures, options) is set across the page in reading order, then the two columns
% resume. The question's head - its number, stem and first figure (\qfullsplit marks the end of the head) -
% is kept on one page, so a number is never left at a page foot with its figure overleaf; the rest may break
% across pages normally, so no big gaps are left. Without a \qfullsplit, the whole question is kept together
% if it is under 45% of a page. Never write these yourself.
\newif\ifqf@full
\newif\ifqf@boxed
\newsavebox{\qf@box}
\newcommand{\qf@open}{\global\qf@fulltrue\par\global\qf@boxedtrue
  \setbox\qf@box\vbox\bgroup\hsize\textwidth\linewidth\textwidth}
\newcommand{\qf@emit}[1]{\par\egroup\global\qf@boxedfalse
  \ifdim\dimexpr\ht\qf@box+\dp\qf@box\relax<#1\textheight\noindent\usebox{\qf@box}\par\else\unvbox\qf@box\fi}
\newcommand{\qfullsplit}{\ifqf@boxed\qf@emit{0.85}\fi}
\newcommand{\qf@close}{\ifqf@boxed\qf@emit{0.45}\fi\global\qf@fullfalse}
\newcommand{\qfullbegin}{\end{multicols}\qf@open}
\newcommand{\qfullend}{\qf@close\begin{multicols}{2}}
\newcommand{\qfullstart}{\qf@open}
\newcommand{\qfulldone}{\qf@close}
% \qfig[scale]{id}: block figure, centred, directly where it is written (size set by the tool).
% Fallback if a wide figure is not inside a full-width question: break out of the columns in place - no floats.
\newcommand{\qfig}[2][1]{\qf@check{#2}%
  \ifqf@full\qf@blk{#1}{#2}\else
  \ifnum\@nameuse{qf@wide@#2}=1 \expandafter\@firstoftwo\else\expandafter\@secondoftwo\fi
  {\end{multicols}\begin{center}\qf@img{#1}{\textwidth}{#2}\end{center}\begin{multicols}{2}}%
  {\qf@blk{#1}{#2}}\fi}
\newcommand{\qf@blk}[2]{\par\smallskip{\centering\qf@img{#1}{\linewidth}{#2}\par}\smallskip}
% \qfigi[scale]{id}: inline figure (inside option lists / tables / minipages); never breaks columns.
\newcommand{\qfigi}[2][1]{\qf@check{#2}\qf@img[valign=c]{#1}{\linewidth}{#2}}
\makeatother
% structure macros - lint reads these to check completeness and figure placement
\newcommand{\papersec}[1]{\par\bigskip{\centering\Large\bfseries PAPER-#1\par}\medskip}
\newcommand{\partsec}[2]{\par\medskip\noindent{\large\bfseries PART-#1: #2}\par\smallskip}
\newcommand{\sect}[1]{\par\smallskip\noindent\textbf{#1}\par}
\newcommand{\qn}[1]{\par\smallskip\noindent\textbf{#1.}~}
\newcommand{\paraq}[2]{\par\smallskip\noindent\textbf{Paragraph for Questions #1 and #2}\par}
\newenvironment{qoptions}{\begin{itemize}[leftmargin=1.8em,itemsep=1pt,topsep=2pt,label={}]}{\end{itemize}}
\newcommand{\opt}[2]{\item[(#1)] #2}
% figure registry (generated by score2_tools.py finalize)
\qfigdef{p1_phy_s1_q01_1}{TOPS - Test 01 (31 Jul 2022)-img/p1_phy_s1_q01_1.png}{55.6}{0}
\qfigdef{p1_phy_s1_q02_1}{TOPS - Test 01 (31 Jul 2022)-img/p1_phy_s1_q02_1.png}{124.2}{0}
\qfigdef{p1_phy_s1_q03_1}{TOPS - Test 01 (31 Jul 2022)-img/p1_phy_s1_q03_1.png}{34.5}{0}
\qfigdef{p1_phy_s1_q04_1}{TOPS - Test 01 (31 Jul 2022)-img/p1_phy_s1_q04_1.png}{134.7}{0}
\qfigdef{p1_phy_s1_q05_1}{TOPS - Test 01 (31 Jul 2022)-img/p1_phy_s1_q05_1.png}{82.0}{0}
\qfigdef{p1_phy_s1_q06_1}{TOPS - Test 01 (31 Jul 2022)-img/p1_phy_s1_q06_1.png}{126.3}{0}
\qfigdef{p1_phy_s1_q07_1}{TOPS - Test 01 (31 Jul 2022)-img/p1_phy_s1_q07_1.png}{97.9}{0}
\qfigdef{p1_phy_s2_para9_10_1}{TOPS - Test 01 (31 Jul 2022)-img/p1_phy_s2_para9_10_1.png}{134.7}{0}
\qfigdef{p1_phy_s2_para11_12_1}{TOPS - Test 01 (31 Jul 2022)-img/p1_phy_s2_para11_12_1.png}{81.8}{0}
\qfigdef{p1_phy_s2_para11_12_2}{TOPS - Test 01 (31 Jul 2022)-img/p1_phy_s2_para11_12_2.png}{102.4}{0}
\qfigdef{p1_phy_s2_para11_12_3}{TOPS - Test 01 (31 Jul 2022)-img/p1_phy_s2_para11_12_3.png}{143.3}{0}
\qfigdef{p1_phy_s3_q01_1}{TOPS - Test 01 (31 Jul 2022)-img/p1_phy_s3_q01_1.png}{127.7}{0}
\qfigdef{p1_phy_s3_q02_1}{TOPS - Test 01 (31 Jul 2022)-img/p1_phy_s3_q02_1.png}{120.0}{0}
\qfigdef{p1_phy_s3_q05_1}{TOPS - Test 01 (31 Jul 2022)-img/p1_phy_s3_q05_1.png}{179.6}{0}
\qfigdef{p1_phy_s3_q06_1}{TOPS - Test 01 (31 Jul 2022)-img/p1_phy_s3_q06_1.png}{119.6}{0}
\qfigdef{p1_chem_s1_q02_A}{TOPS - Test 01 (31 Jul 2022)-img/p1_chem_s1_q02_A.png}{108.0}{0}
\qfigdef{p1_chem_s1_q05_1}{TOPS - Test 01 (31 Jul 2022)-img/p1_chem_s1_q05_1.png}{73.4}{0}
\qfigdef{p1_chem_s1_q05_2}{TOPS - Test 01 (31 Jul 2022)-img/p1_chem_s1_q05_2.png}{84.0}{0}
\qfigdef{p1_chem_s1_q06_1}{TOPS - Test 01 (31 Jul 2022)-img/p1_chem_s1_q06_1.png}{81.3}{0}
\qfigdef{p1_chem_s2_para9_10_1}{TOPS - Test 01 (31 Jul 2022)-img/p1_chem_s2_para9_10_1.png}{92.6}{0}
\raggedcolumns
\makeatletter\@beginparpenalty=10000\makeatother
\newcolumntype{L}[1]{>{\raggedright\arraybackslash}p{#1}}
\begin{document}
\begin{center}{\Large\bfseries TOPS --- Test 01 (31 Jul 2022)}\end{center}
\begin{multicols}{2}
\raggedcolumns
\papersec{1}
\partsec{1}{PHYSICS}
\sect{SECTION-I(i)}
\qn{1} In the figure shown, mass of $A$ is 5 kg. All pulleys and strings are ideal, then :
\qfig{p1_phy_s1_q01_1}
\begin{qoptions}
\opt{A}{If acceleration of $A$ is 1 m/s$^2$ upwards, then acceleration of $B$ is 3 m/s$^2$ downwards.}
\opt{B}{If acceleration of $A$ is 3 m/s$^2$ upwards, then acceleration of $B$ is 1 m/s$^2$ downwards.}
\opt{C}{If mass of $B$ is 20 kg, then acceleration of $B$ is $\dfrac{3g}{13}$ downwards.}
\opt{D}{If mass of $B$ is 20 kg, then acceleration of $A$ is $\dfrac{g}{13}$ upwards.}
\end{qoptions}
\qn{2} In the figure show, the man pulls the strings with velocity $u$ from both ends and the block of mass $M$ always remains horizontal. Then,
\qfig{p1_phy_s1_q02_1}
\begin{qoptions}
\opt{A}{Velocity of block of mass $M$ is $\dfrac{3u}{4}$.}
\opt{B}{Velocity of pulleys $P_1$ is $2u$.}
\opt{C}{Velocity of pulley $P_2$ is $3u/4$.}
\opt{D}{Velocity of pulley $P_3$ is $u/2$.}
\end{qoptions}
\qn{3} A block is attached to two identical springs as shown in figure. The extension in upper spring is twice the compression in lower spring. Then :
\qfig{p1_phy_s1_q03_1}
\begin{qoptions}
\opt{A}{If upper spring is cut, just after that the block will accelerate upward with $\dfrac{2g}{3}$.}
\opt{B}{If upper spring is cut, just after that the block will accelerate downward with $\dfrac{2g}{3}$.}
\opt{C}{If lower spring is cut, just after that the block will accelerate upward with $\dfrac{g}{3}$.}
\opt{D}{If lower spring is cut, just after that the block will accelerate downward with $\dfrac{g}{3}$.}
\end{qoptions}
\qn{4} Two blocks of masses $m_1$ and $m_2$ are connected with massless undeformed spring and placed over a long plank moving with an acceleration `$a$' as shown in figure. The coefficient of friction between the blocks and platform is `$\mu$'.
\qfig{p1_phy_s1_q04_1}
\begin{qoptions}
\opt{A}{Spring will be stretched if $a>\mu g$}
\opt{B}{Spring will be compressed if $a\le\mu g$}
\opt{C}{Spring will neither be compressed nor be stretched only if $a\le\mu g$}
\opt{D}{Spring will be in its natural length under all conditions}
\end{qoptions}
\qn{5} A plank of mass $m_1=8$ kg with a bar of mass $m_2=2$ kg placed on its rough surface, lie on a smooth floor of elevator ascending with an acceleration $g/4$. The coefficient of friction is $\mu=1/5$ between $m_1$ and $m_2$. A horizontal force $F=30$ N is applied to the plank. Then, in the reference frame of elevator ($g=10$ m/s$^2$)
\qfig{p1_phy_s1_q05_1}
\begin{qoptions}
\opt{A}{The acceleration of bar is 2.5 m/s$^2$}
\opt{B}{The acceleration of bar is 3.5 m/s$^2$}
\opt{C}{The acceleration of plank is 5 m/s$^2$}
\opt{D}{The acceleration of plank is $\dfrac{25}{8}$ m/s$^2$.}
\end{qoptions}
\qn{6} Two blocks with masses $M_1$ and $M_2$ of 10 kg and 20 kg respectively are placed as shown in figure coefficient of friction $=0.2$ between all surfaces, then just after the system is released from rest ($g=10$ m/s$^2$) string and pulley are light. Neglect the rotation of $M$.
\qfig{p1_phy_s1_q06_1}
\begin{qoptions}
\opt{A}{Tension in string is 306 N nearly.}
\opt{B}{Tension in string is 250 N nearly.}
\opt{C}{Acceleration of $M_2$ block is 4.7 m/s$^2$ nearly.}
\opt{D}{Acceleration of $M_2$ block is 6.5 m/s$^2$ nearly.}
\end{qoptions}
\qn{7} A solid cube of mass 5 kg is placed on a rough horizontal surface, in $xy$--plane as shown. The friction coefficient between the surface and the cube is 0.4. An external force $\vec F=\left(12\hat i+16\hat j+20\hat k\right)$ N is applied on the cube. (use $g=10$ m/s$^2$)
\qfig{p1_phy_s1_q07_1}
\begin{qoptions}
\opt{A}{The block starts slipping over the surface}
\opt{B}{The friction force on the cube by the surface is 10 N.}
\opt{C}{The friction force acts in $xy$--plane at angle $127^\circ$ with the positive $x$--axis in clockwise direction.}
\opt{D}{The contact force exerted by the surface on the cube is $10\sqrt{10}$ N.}
\end{qoptions}
\qn{8} A uniform rod of length `$L$' and mass `$M$' has been placed on a rough horizontal surface. The horizontal force `$F$' applied on the rod at one end such that the rod is just in the state of rest and about to slip. If the coefficient of friction varies according to the relation $\mu=Kx$. (where $x$ is the distance from the opposite end). Where `$K$' is $+$ve constant. Then :
\begin{qoptions}
\opt{A}{Tension at mid point is $F/3$}
\opt{B}{Tension at mid point is $F/4$.}
\opt{C}{Tension at $x=\dfrac L4$ is $\dfrac F{16}$}
\opt{D}{Tension at $x=\dfrac{3L}4$ is $\dfrac{9F}{16}$}
\end{qoptions}
\sect{SECTION-I(ii)}
\paraq{9}{10} In the arrangement shown in figure the masses $m$ of the bar and $M$ of the wedge, as well as the wedge angle $\alpha$, are known. The masses of the pulley and the thread are negligible. The friction is absent.
\qfig{p1_phy_s2_para9_10_1}
\qn{9} Which of the following statement(s) is/are correct ?
\begin{qoptions}
\opt{A}{The acceleration of wedge $M$ is $\dfrac{mg\sin\alpha}{M+m(1-\cos\alpha)}$}
\opt{B}{The acceleration of wedge $M$ is $\dfrac{mg\sin\alpha}{M+2m(1+\cos\alpha)}$}
\opt{C}{The acceleration of bar relative to wedge is $\dfrac{mg\sin\alpha}{M+2m(1-\cos\alpha)}$}
\opt{D}{The acceleration of bar relative to wedge is $\dfrac{mg\sin\alpha}{M+2m(1+\cos\alpha)}$}
\end{qoptions}
\qn{10} If in the previous situation ground were rough to prevent slipping while other surfaces still smooth, then which of the following options is incorrect ?
\begin{qoptions}
\opt{A}{The frictional force between wedge and ground is $mg\sin\alpha$.}
\opt{B}{The tension in the string is $mg\sin\alpha$.}
\opt{C}{Normal force acting on the wedge by the ground is $(M+m)g$}
\opt{D}{Normal force acting on the bar by the wedge is $mg$.}
\end{qoptions}
\paraq{11}{12} When two rigid bodies move over each other, their intersection point might not be moving with the velocity of either of them. For example, consider two rods $AB$ and $CD$, connected by pin $P$ that moves in slots made in the rods.
\qfig{p1_phy_s2_para11_12_1}
If the rod $AB$ is moving with velocity $v_1$ and rod $CD$ is moving with velocity $v_2$, their intersection point (pin $P$) moves as shown:
\qfig{p1_phy_s2_para11_12_2}
Where $A'B'$ and $C'D'$ are the new positions of the rods, while $P'$ is the new position of the pin $P$. Here the pin $P$ is having a velocity that is the sum of the velocities of the two rods. But in more complicated shapes of the sliders, this direct approach might not work, and calculus is required to calculate the velocity or acceleration of the intersection point.

In the given diagram, there is an elliptical slot as shown. The equation of the ellipse is $\dfrac{x^2}{4}+y^2=1$, where $x$ and $y$ are the coordinates in same unit. There is a slider with two pins $A$ and $B$, which are free to move along the slider so that they always remain on the elliptical slot. The slider is moved horizontally with constant speed 10 m/s. Answer the following questions.
\qfig{p1_phy_s2_para11_12_3}
\qn{11} When $x$-coordinate of the slider is 1 m, the speed of Pin-$B$ is?
\begin{qoptions}
\opt{A}{Zero}
\opt{B}{10 m/s}
\opt{C}{2.9 m/s}
\opt{D}{10.4 m/s}
\end{qoptions}
\qn{12} When $x$-coordinate of the slider is 1 m, the magnitude of acceleration of Pin-$B$ is?
\begin{qoptions}
\opt{A}{Zero}
\opt{B}{38.4 cm/s$^2$}
\opt{C}{38.5 m/s$^2$}
\opt{D}{42 cm/s$^2$}
\end{qoptions}
\sect{SECTION-II}
\qn{1} In the arrangement shown in figure, the block of mass $m=2$ kg lies on the wedge of mass $M=8$ kg. The surfaces are smooth and pulley \& strings are massless. If the initial acceleration of wedge is $\dfrac{30\sqrt a}{18+b}$ m/s$^2$, then find $a+b$.
\qfig{p1_phy_s3_q01_1}
\qn{2} The maximum value of mass of block $C$ (in kg) so that neither $A$ nor $B$ moves is (Given that mass of $A$ is 100 kg and that of $B$ is 140 kg. Pulleys are smooth and friction coefficient between $A$ and $B$ and between $B$ and horizontal surface is ($\mu=0.3$) $g=10$ m/s$^2$)
\qfig{p1_phy_s3_q02_1}
\qn{3} The minimum number of unequal, non-coplanar forces acting on an object so that it remains at rest is :
\qn{4} A heavy body of mass 25 kg is to be dragged along a horizontal plane ($\mu=1/\sqrt3$). If the least force required to move the body is $25x$ newton, find $x$.
\qn{5} Initial velocity of the block $\bar v=45\hat i$ while external force on it is $\bar F=25\left(-\hat i\right)$. It coefficient of static and kinetic friction are 0.3 and 0.2 respectively then distance travelled by the block in 12 sec is $25x$ m then $x$ value is ($g=10$ m/s$^2$) :
\qfig{p1_phy_s3_q05_1}
\qn{6} A block $A$ of mass $M$ rests on a wedge $B$ of mass $2M$ of inclination $\theta$. There is sufficient friction between $A$ and $B$ so that $A$ does not slip on $B$. If there is no friction between $B$ and ground, the compression in spring is
\qfig{p1_phy_s3_q06_1}
\partsec{2}{CHEMISTRY}
\sect{SECTION-I(i)}
\qn{1} In which pair both the species are non existing?
\begin{qoptions}
\opt{A}{SH$_4$ ; SH$_6$}
\opt{B}{OF$_4$ ; OF$_6$}
\opt{C}{SF$_4$ ; OF$_4$}
\opt{D}{NOF$_3$ ; NF$_5$}
\end{qoptions}
\qn{2} Which of the following ground state electronic configuration(s) is/are correct?
\begin{qoptions}
\opt{A}{\qfigi{p1_chem_s1_q02_A}}
\opt{B}{$\mathrm{[Ar]}\left(3d^1_{xy}=3d^1_{yz}\right)4s^2$}
\opt{C}{$\mathrm{[Ne]}3s^2\left(3p_y^2=3p_z^1=3p_x^0\right)$}
\opt{D}{$\mathrm{[Rn]}\left(6d^1_{x^2-y^2}=6d^1_{z^2}\right)7s^2$}
\end{qoptions}
\qn{3} Consider the following statements :

I. Enthalpy of formation is always calculated from its constituent elements in their solid state.

II. Hydration energy of ion is always amount of energy released.

III. Lattice energy is amount of energy supplied to an ionic solid so that it is melted into constituent ions
\begin{qoptions}
\opt{A}{I \& II are correct}
\opt{B}{II is correct}
\opt{C}{I \& III are incorrect}
\opt{D}{I, II \& III are incorrect}
\end{qoptions}
\qn{4} Which of the following molecules will show at least one angle $90^\circ$ if repulsion due to lone pair is not considered.
\begin{qoptions}
\opt{A}{SeCl$_4$}
\opt{B}{XeF$_2$}
\opt{C}{XeF$_5^-$}
\opt{D}{BrF$_5$}
\end{qoptions}
\qn{5} Consider the following two structures :
\qfig{p1_chem_s1_q05_1}
\qfig{p1_chem_s1_q05_2}
Which is/are correct?
\begin{qoptions}
\opt{A}{\% s character in P--H bond is more in II as compared to I}
\opt{B}{Both have p$_\pi$-d$_\pi$ bonds}
\opt{C}{Oxidation number of phosphorus is higher in (II)}
\opt{D}{Both have equal number of acidic hydrogen}
\end{qoptions}
\qn{6} Consider the skeleton structure of thiocyanic acid.
\qfig{p1_chem_s1_q06_1}
Which is/are correct regarding above molecule.
\begin{qoptions}
\opt{A}{(Bond order)$_{\mathrm I}>$ (Bond order)$_{\mathrm{II}}$}
\opt{B}{(Bond order)$_{\mathrm{II}}>$ (Bond order)$_{\mathrm I}$}
\opt{C}{Geometry of sulphur is same as N in HNCS}
\opt{D}{Bond angle around carbon is same as in HNCS}
\end{qoptions}
\qn{7} Which of the following statement(s) is/are correct?
\begin{qoptions}
\opt{A}{All C--N bond lengths in C$_3$N$_3$(N$_3$)$_3$ are identical}
\opt{B}{AsH$_3$ has less $\mathrm{H-\hat X-H}$ bond angle compared to $\mathrm{H-\hat X-H}$ bond angle in NH$_3$ ($X\to$ central atom)}
\opt{C}{CH$_3$Cl has more dipole moment than CH$_3$F}
\opt{D}{C$_3$H$_4$ has both $\pi$-bonds in different planes}
\end{qoptions}
\qn{8} Generally the 5d series elements show similar radius with their 4d series counterparts. This is supported by the fact that
\begin{qoptions}
\opt{A}{Due to Lanthanide contraction, the 4f subshell experience poor shielding by inner electrons}
\opt{B}{Due to Lanthanide contraction, the 5d \& 6s electrons are shielded poorly by 4f subshell.}
\opt{C}{Along with 4f, the inner `d' subshells (3d \& 4d) also show poor shielding making the size decrease to a greater extent}
\opt{D}{As number of electrons in `s' subshell increases, size decreases}
\end{qoptions}
\columnbreak
\sect{SECTION-I(ii)}
\paraq{9}{10} Consider the following structure.
\qfig{p1_chem_s2_para9_10_1}
\qn{9} If the above structure is P$_4$O$_{10}$ then :-
\begin{qoptions}
\opt{A}{$x>y$}
\opt{B}{$\theta_1>\theta_2$}
\opt{C}{$x=y$}
\opt{D}{$\theta_1<\theta_2$}
\end{qoptions}
\qn{10} If the above structure is P$_4$S$_{10}$, then
\begin{qoptions}
\opt{A}{Total four 2p$_\pi$-3d$_\pi$ bonds are present}
\opt{B}{$(x)_{\mathrm{P_4O_{10}}}>(x)_{\mathrm{P_4S_{10}}}$}
\opt{C}{Oxidation state of P in P$_4$S$_{10}$ is more compared to P in P$_4$O$_{10}$}
\opt{D}{Have 16 $\sigma$ \& 4 $\pi$ bonds}
\end{qoptions}
\paraq{11}{12} An electron of an element ``X'' having $m=+1$ is involved in bonding with an electron having $m=+1$ of another element Y. The bond formed by the electrons is not formed by head-on overlapping of orbitals.
\qn{11} Which of the following specie can have a bond mentioned above ($X\to$ central atom ; $Y\to$ surrounding atom)
\begin{qoptions}
\opt{A}{PH$_3$}
\opt{B}{COCl$_2$}
\opt{C}{SF$_4$}
\opt{D}{PCl$_6^-$}
\end{qoptions}
\qn{12} If the electron in Y is having $m=+1$, \& X has electron having $m=+2$, and bond formed is not formed by head on overlapping then specie that cannot have such a bond is ($X\to$ central atom ; $Y\to$ surrounding atom).
\begin{qoptions}
\opt{A}{PO$_4^{3-}$}
\opt{B}{NO$_3^-$}
\opt{C}{SO$_2$}
\opt{D}{S$_3$O$_9$}
\end{qoptions}
\sect{SECTION-II}
\qn{1} Number of species that have a coordinate bond among the following.

KCN, NH$_4$SO$_3$F, FXeOSO$_2$Cl, PCl$_6^-$, HNC, HONO$_2$, CO, NO$_2$Cl
\qn{2} Percentage ionic character of a heterodiatomic molecule (XY) is 19.5\%.

``Y'' exists as Y$_2$ in nature \& bond length of Y$_2$ is 200 pm. Radius of X is 80 pm. Then the value of bond length of (X--Y) (in pm) is \underline{\hspace{1.5cm}}.
\qn{3} Element ``A'' only forms acidic oxides. ``A'' belongs to the period containing Bridging elements.

If A exists in nature as A$_2$ ; then the difference in oxidation states of A in most acidic \& least acidic oxides of A is \underline{\hspace{1cm}}.
\qn{4} Number of species that have identical structures both according to Lewis theory \& according to maximum covalency among the following.

PO$_4^{3-}$, H$_2$SO$_4$, NO$_2^-$, COCl$_2$, C$_3$O$_2$, CO$_3^{2-}$, NH$_4^+$, HPO$_3^{2-}$
\qn{5} Consider the orbitals

$d_{xy}$, $s$, $d_{yz}$, $d_{z^2}$, $d_{x^2-y^2}$, $p_y$, $p_x$, $p_z$

If the number of orbitals among the following that have zero nodal planes are $X$ and orbitals that can form a delta bond are $Y$, then $(Y-X)$ is.
\qn{6} Number of correct order(s) among the following

(i) Se $>$ O (EA)

(ii) Na $>$ Al (EA)

(iii) H$_2$O $>$ NH$_3$ (Bond Angle)

(iv) Pd $>$ Ni (IE)

(v) N$_3^-$ $>$ CO$_2$ (number of electrons)

(vi) Mn $>$ Fe (IE$_3$)

(vii) B $>$ H (EN)
\partsec{3}{MATHEMATICS}
\sect{SECTION-I(i)}
\qn{1} Let $a,b\in\mathbb Q$ such that $\dfrac{14\sqrt3+17}{4-\sqrt3}=a\sqrt3+b$ then which of the following is true ?
\begin{qoptions}
\opt{A}{$\dfrac ab$ is an integer}
\opt{B}{$\sqrt{\dfrac ab}$ is irrational number}
\opt{C}{$\sqrt{b-a}$ is an irrational number}
\opt{D}{$\sqrt{a+b}$ is rational number}
\end{qoptions}
\qn{2} Sum of the product of roots of the equation $(x-1)(x-2)\left(x-2^2\right)\left(x-2^3\right)\ldots..\left(x-2^9\right)=0$ taken two at a time is -
\begin{qoptions}
\opt{A}{$\dfrac{2^{20}}3+\dfrac23$}
\opt{B}{$\dfrac{2^{20}}3-\dfrac{3071}3$}
\opt{C}{$\dfrac{2^{20}}3-\dfrac13$}
\opt{D}{$\dfrac{2^{20}}3-\dfrac{3070}3$}
\end{qoptions}
\qn{3} If both roots of $x^2-42x+\lambda=0$ are prime numbers, then possible values of $\lambda$ is/are -
\begin{qoptions}
\opt{A}{185}
\opt{B}{341}
\opt{C}{377}
\opt{D}{437}
\end{qoptions}
\qn{4} If $\alpha$ ($\alpha>0$) is a root of $x^{11}-x^{10}+x^8-x^7+x^5-x^4+x^2-x-20=0$ then $\alpha^{12}$ is greater than
\begin{qoptions}
\opt{A}{59}
\opt{B}{61}
\opt{C}{58}
\opt{D}{57}
\end{qoptions}
\qn{5} If $\dfrac x2$, $|x+1|$, $|x-1|$ are the first three terms of an arithmetic progression then
\begin{qoptions}
\opt{A}{Sum of first 50 terms of the AP can be 970}
\opt{B}{50$^{\text{th}}$ term of the AP can be 40}
\opt{C}{Sum of first 50 terms of an AP can be 2400}
\opt{D}{50$^{\text{th}}$ terms of the AP can be 97}
\end{qoptions}
\qn{6} Which of the following hold(s) good ?
\begin{qoptions}
\opt{A}{$(\log_35)(\log_34)(\log_{10}11)>1$}
\opt{B}{$(\log_23)(\log_{1/3}5)(\log_46)<1$}
\opt{C}{$\log_{100}7>\log_{29}3$}
\opt{D}{$\log_\pi5>\log_e5$}
\end{qoptions}
\qn{7} If $5\sec x-3\tan x=6$, then $6\sin x-3\cos x$ can be -
\begin{qoptions}
\opt{A}{$\sqrt5$}
\opt{B}{$2\sqrt5$}
\opt{C}{$-2\sqrt5$}
\opt{D}{$-\sqrt5$}
\end{qoptions}
\qn{8} If $x\ne\dfrac{n\pi}2$ ($n$ is integer), then the expression $\dfrac{\sin^3x}{1+\cos x}+\dfrac{\cos^3x}{1-\sin x}$ is equivalent to
\begin{qoptions}
\opt{A}{$\sqrt2\cos\left(\dfrac\pi4-x\right)$}
\opt{B}{$\sqrt2\cos\left(\dfrac\pi4+x\right)$}
\opt{C}{$\sqrt2\sin\left(\dfrac\pi4-x\right)$}
\opt{D}{$\sqrt2\sin\left(\dfrac\pi4+x\right)$}
\end{qoptions}
\sect{SECTION-I(ii)}
\paraq{9}{10} Consider the polynomial equation $\displaystyle\prod_{k=1}^{n}\left(x^2-28x+k\right)=0$.
\qn{9} If $n=196$, the product of all real \& distinct roots of the above polynomial equation is-
\begin{qoptions}
\opt{A}{$\displaystyle\prod_{r=1}^{196}(r)$}
\opt{B}{$\displaystyle\prod_{r=1}^{195}(r)$}
\opt{C}{$14\cdot\displaystyle\prod_{r=1}^{195}(r)$}
\opt{D}{$14\cdot\displaystyle\prod_{r=1}^{196}(r)$}
\end{qoptions}
\qn{10} If $n=100$, the sum of all real roots of the above polynomial equation is-
\begin{qoptions}
\opt{A}{2800}
\opt{B}{2786}
\opt{C}{5474}
\opt{D}{5488}
\end{qoptions}
\paraq{11}{12} Consider the series
\par\noindent\hfill$\begin{aligned}&\underbrace{1^2x^1}_{\text{1 term}}+\underbrace{2^2x^2+2^2x^3}_{\text{2 terms}}\\&+\underbrace{4^2x^4+4^2x^5+4^2x^6+4^2x^7}_{\text{4 terms}}\\&+\underbrace{8^2x^8+8^2x^9+\cdots}_{\text{8 terms}}+\cdots\end{aligned}$\hfill\null\par
\qn{11} If $x=1$, then sum of first 1023 terms is equal to
\begin{qoptions}
\opt{A}{$\dfrac{4^{20}-1}3$}
\opt{B}{$\dfrac{8^{10}-1}7$}
\opt{C}{$\left(2^{1023}-1\right)$}
\opt{D}{$16^5$}
\end{qoptions}
\qn{12} If $x=-1$, then sum of first 1023 terms is equal to
\begin{qoptions}
\opt{A}{$-1$}
\opt{B}{$\dfrac{2^{20}-1}3$}
\opt{C}{$\dfrac{4^{10}-1}3$}
\opt{D}{$\dfrac{8^{10}-1}7$}
\end{qoptions}
\sect{SECTION-II}
\qn{1} If $\sqrt{\alpha_1-1}+2\sqrt{\alpha_2-4}=\dfrac{\alpha_1+\alpha_2}2$, where $\alpha_1,\alpha_2\in\mathbb N$, then find the value of $\alpha_1.\alpha_2$
\qn{2} The arithmetic mean of solutions of equation $\sqrt[3]{\sin x-1}+\sqrt[3]{\sin x}+\sqrt[3]{\sin x+1}=0$, lying in $[0,2\pi]$, is
\qn{3} The value of
\par\noindent\hfill$\begin{aligned}&\sqrt{(\log_xy)\left(\log_x\left(x^2y\right)\right)\left(\log_x\left(x^4y\right)\right)\left(\log_x\left(x^6y\right)\right)+16}\\&-\left(\log_x\left(x^2y\right)\right)\left(\log_x\left(x^4y\right)\right)\end{aligned}$\hfill\null\par
$x>1$, $y>1$; is
\qn{4} Find the largest common element of two sets $\{1,4,7,10,\ldots,541\}$ and $\{3,8,13,18,\ldots,498\}$
\qn{5} If $9(27)^{x^2+y+z}+9(27)^{x+y^2+z}+9(27)^{x+y+z^2}=1$, then number of possible real triplet $(x,y,z)$ is
\qn{6} If $x=\dfrac{\sin p}{\cos^2p}$, $y=\dfrac{\cos p}{\sin^2p}$ and $\sin p+\cos p=\dfrac12$ then the value of $x+y$ is
\end{multicols}
\end{document}
